\documentclass[modern,linenumbers]{aastex701}

\usepackage{graphicx}
\usepackage{amsmath}
\usepackage{xcolor}
\usepackage{longtable}

% color macros for editorial notes
\newcommand{\JM}[1]{{\color{blue}#1}}
\newcommand{\todo}[1]{{\color{red}\textbf{TODO:} #1}}

% wrapped-text deluxetable column (deluxetable's preamble scanner chokes
% on a bare p{...} column, so alias it to a single letter)
\newcolumntype{W}{>{\raggedright\arraybackslash}p{3.5in}}

\graphicspath{{./}{figures/}}

\submitjournal{\aj}

\shorttitle{SGA-2025}
\shortauthors{Moustakas et al.}

\begin{document}

\title{The Siena Galaxy Atlas 2025}

\correspondingauthor{John Moustakas}

\author[0000-0002-2733-4559]{John Moustakas}
\email{jmoustakas@siena.edu}
\affiliation{Department of Physics and Astronomy, Siena University,
  515 Loudon Road, Loudonville, NY 12211, USA}

% \author{...}
% \affiliation{...}

\author{Imaging Team}
\email{team@teamtastic.edu}
\affiliation{INSTITUTION, ADDRESS}


\begin{abstract}

We present the Siena Galaxy Atlas 2025 (SGA-2025), a photometric
catalog of $470,625$ large, angularly resolved galaxies selected from
the DESI Legacy Imaging Surveys Data Release 11 (DR11).  The SGA-2025
covers $\sim\!30{,}000$~deg$^{2}$ of the extragalactic sky and
delivers multiwavelength imaging mosaics and ellipse photometry in up
to ten bands: optical $grz$ or $griz$ imaging from DECam (south) and
BASS+MzLS (north), mid-infrared $W1$--$W4$ from unWISE, and
ultraviolet FUV and NUV from GALEX.  For each galaxy we provide
curve-of-growth total magnitudes, the isophotal diameter at a
surface-brightness threshold of $\mu = 26~\mathrm{mag~arcsec^{-2}}$
($D_{26}$), the half-light semi-major axis, isophotal sizes at several
other thresholds, and fitted ellipse geometry.  We cross-match
spectroscopic redshifts from DESI DR1, the NASA/IPAC Extragalactic
Database (NED), and the Local Volume Database (LVD) and include them
in the catalog.  The SGA-2025 supersedes the SGA-2020
\citep{moustakas23a}, more than doubling the number of catalogued
galaxies while extending coverage into the ultraviolet and
mid-infrared.  We make the full catalog, mosaics, and companion
analysis code publicly available at
\url{https://sga.legacysurvey.org}.

\end{abstract}

\keywords{catalogs --- galaxies: photometry --- galaxies: fundamental
parameters --- surveys}

%%----------------------------------------------------------------------
\section{Introduction}
\label{sec:intro}
%%----------------------------------------------------------------------

The Siena Galaxy Atlas (SGA) pursues a goal distinct from that of most
galaxy catalogs: rather than selecting objects by apparent magnitude or
spectroscopic redshift, we assemble the most complete
angular-diameter-limited sample of large, resolved galaxies ever built,
spanning the full extragalactic sky accessible to modern digital
imaging.  This selection function matters because spatially resolved
photometry of nearby, angularly large galaxies---precise sizes, shapes,
and fluxes measured homogeneously across the sky---underpins a wide
range of extragalactic astrophysics: the galaxy luminosity and
stellar-mass functions, the size--luminosity and Tully--Fisher
relations, group and environmental studies, and the selection of
targets and training sets for spectroscopic surveys and
machine-learning morphology classifiers.  Only galaxies near or large
enough to be spatially well resolved reveal the internal structure,
bars, rings, disks, and low-surface-brightness features invisible in
unresolved photometry, that connects present-day galaxy morphology to
the star formation, merger, and feedback histories that build it over
cosmic time (see \citealt{moustakas23a} for a full discussion of the
motivating science questions).  Yet these same large, angularly
extended galaxies are the ones most poorly served by standard
photometric pipelines, which are tuned to the point-like and marginally
resolved sources that dominate magnitude-limited catalogs: their
extended light profiles are prone to shredding into multiple spurious
detections, and the sky background beneath them is easily
over-subtracted, so that even deep, wide-area digital surveys with
excellent point-source photometry can remain surprisingly incomplete
for precisely the galaxies that carry the most spatial information per
object.  We designed the SGA from the outset to correct for these
failure modes directly, rather than to inherit them from a
general-purpose source-detection pipeline.

This structural information has a long observational heritage, from
Charles Messier's eighteenth-century catalog of nebulae through the New
General and Index Catalogues and, eventually, the Third Reference Catalogue
of Bright Galaxies \citep[RC3;][]{de-vaucouleurs91a}, and, more recently, a
generation of wide-area, multiwavelength atlases in the near-infrared,
ultraviolet, and optical (see \citealt{moustakas23a} and references therein
for the full lineage, and \citealt{christie24a} for a recent review of the
broader landscape of nearby-galaxy surveys).  With the exception of the
NASA-Sloan Atlas, however, these atlases' optical sizes, shapes, and total
magnitudes still trace back to the RC3's photographic-plate measurements.
No comprehensive, all-sky photometric catalog of resolved galaxies built
entirely from modern, wide-field digital imaging existed until the Siena
Galaxy Atlas 2020 (SGA-2020; \citealt{moustakas23a}), a $grz$+$W1$--$W4$
atlas of $383{,}620$ galaxies over the $\approx\!20{,}000$~deg$^2$ DESI
Legacy Imaging Surveys Data Release 9 footprint.  Since the SGA-2020's
publication, several new value-added galaxy compilations have extended
this landscape further---among them REGALADE \citep{tranin26a} and
HECATEv2 \citep{kyritsis26a}, both aimed principally at time-domain and
multi-messenger applications (below)---but these are volume- or
magnitude-limited to a few hundred megaparsecs rather than
angular-diameter-selected, and to our knowledge none has yet matched the
SGA's combination of all-sky coverage, multiwavelength depth, and
angular-diameter completeness for large, resolved galaxies.

The DESI Legacy Imaging Surveys have since concluded observations with
Data Release 11 (DR11; \citealt{dey19a}), delivering deeper, wider, and
better-calibrated $grz$/$griz$ imaging than any previous data release,
together with forced photometry from the Tractor source-fitting code
\citep{lang16a}.  This imaging underlies the DESI spectroscopic survey
itself, whose Bright Galaxy Survey needs a high-quality photometric
catalog of large-diameter galaxies to guard against the shredding and
spectroscopic incompleteness that such objects otherwise cause.  We
combine the DR11 optical imaging with full-depth unWISE mid-infrared
coadds \citep{lang14a, meisner21a} and the GALEX ultraviolet sky
\citep{martin05a} to assemble large-galaxy photometry spanning four
decades in wavelength, from the far-ultraviolet to $22~\mu$m.

In the years since its release, the SGA-2020 has already served as the
foundation for a broad range of extragalactic science, underscoring the
value of a systematic, angular-diameter-limited galaxy sample as a
community resource rather than a single-purpose product.  The DESI
Peculiar Velocity Survey selects its Tully--Fisher and Fundamental-Plane
targets directly from SGA-2020 morphologies and sizes
\citep{saulder23a}, placing spectroscopic fibers on the centers and
major axes of spiral galaxies identified in the SGA-2020 to build its
Tully--Fisher distance catalog \citep{douglass26a}.  Visual and
machine-assisted searches of SGA-2020
cutouts have produced systematic censuses of low-surface-brightness
tidal features around tens of thousands of galaxies \citep{sola25a},
while its all-sky coverage has enabled studies of the alignment between
galaxy spin and the cosmic web using, in the words of one such study,
``the largest catalogue of galaxies publicly available''
\citep{muralichandran25a}.  Cosmological simulations, in turn, use the
SGA-2020's homogeneous sizes and colors as an observational benchmark
against which to validate predicted galaxy structural properties
\citep{baes24a, baes24b}.  We expect the SGA-2025's larger sample size, extended
ultraviolet and mid-infrared coverage, and improved completeness to
support an even wider range of such applications.

A large, angular-diameter-selected sample is also, somewhat
serendipitously, well matched to the demands of time-domain and
multi-messenger astrophysics: because apparent angular size favors
intrinsically large, massive, and relatively nearby galaxies, the SGA
preferentially captures the galaxy population most likely to host
gravitational-wave mergers, supernovae, and other transients.  This
match is becoming more consequential as the scale of transient
discovery grows; the Vera C.\ Rubin Observatory's Legacy Survey of
Space and Time is expected to generate on the order of $10^7$ alerts
per night, and real-time association of each transient with its host
galaxy is central to triaging that stream, prioritizing follow-up, and
flagging the rare events---kilonovae, tidal disruption events, fast
radio burst counterparts---that justify limited spectroscopic and
photometric resources \citep{andreoni24a}.  The reference catalogs that
presently support this work, such as GLADE+ \citep{dalya22a} and the
NED Local Volume Sample \citep{cook23a}, remain complete only to a few
hundred megaparsecs in luminosity-weighted galaxy light, and forecasts
for the fourth LIGO--Virgo--KAGRA observing run project that Rubin
alone will increase the annual rate of electromagnetic-counterpart
detections by roughly a factor of five relative to the Zwicky Transient
Facility \citep{kiendrebeogo23a}---a gain that depends directly on
having a sufficiently deep and complete galaxy catalog against which to
search.  The same logic extends to localizing the hosts of fast radio
bursts \citep{law24a} and to galaxy-catalog-based dark-siren cosmology
\citep{abbott23a}, both of which trade directly on catalog
completeness and depth.  We designed the SGA-2025 with these
applications explicitly in mind: the parent-sample improvements we
describe below (Section~\ref{sec:parent}) close exactly the
completeness gaps that limit a galaxy catalog's value for real-time
transient follow-up, while its uniform, multiwavelength photometry
supports the host-galaxy characterization that such follow-up requires.

Here we present the SGA-2025, which exploits this improved imaging to
deliver mosaics and ellipse photometry for $\sim\!486{,}000$ galaxies
across the DR11 extragalactic footprint.  Concluding the SGA-2020, we
identified its principal limitation as the heterogeneity and
incompleteness it inherited from relying on HyperLeda \citep{makarov14a}
as its near-exclusive parent-catalog input, and we proposed remedying it
by detecting large galaxies directly rather than depending solely on
external compilations.  The SGA-2025 delivers on that plan: rather than
starting from the SGA-2020 itself, we rebuild the parent sample from
scratch, merging HyperLeda with five additional source lists---including
the SGA-2020, restricted to the DR11 footprint, as one input among several
(Section~\ref{sec:parent})---so that the SGA-2025 recovers the great
majority of SGA-2020 galaxies while adding many newly identified ones and,
following fresh visual inspection, occasionally dropping a previous entry
that does not survive the SGA-2025's own selection cuts.

Beyond photometry, we substantially expand the SGA's ancillary
information relative to the 2020 release.  We cross-match spectroscopic
redshifts from the Local Volume Database \citep[LVD;][]{pace25a}, DESI
DR1, SDSS DR17, and NED; adopt distances from a priority-ranked
compilation of high-precision literature measurements, LVD, and NED,
falling back to the Hubble flow when no direct measurement exists; and
train a photometric-redshift model on the resulting spectroscopic
subsample (Section~\ref{sec:catalog}).

We organize the remainder of the paper as follows.
In Section~\ref{sec:imaging} we describe the sky footprint and the
optical, mid-infrared, and ultraviolet imaging data.
In Section~\ref{sec:parent} we describe the input catalog, sample
selection, and group catalog.
In Section~\ref{sec:pipeline} we describe the photometric pipeline,
including image coaddition, Tractor source detection, and ellipse fitting.
In Section~\ref{sec:catalog} we describe the construction, structure, and
data model of the released catalog, together with our spectroscopic
redshift cross-matching and adopted-distance compilation.
We summarize in Section~\ref{sec:summary}.

Throughout we adopt a flat $\Lambda$CDM cosmology with
$H_0 = 70~\mathrm{km~s^{-1}~Mpc^{-1}}$ and $\Omega_m = 0.3$.
All magnitudes are in the AB system.

%%----------------------------------------------------------------------
\section{Parent Sample and Group Catalog}
\label{sec:parent}
%%----------------------------------------------------------------------

\subsection{Building the Parent Sample}
\label{sec:parent:summary}

We assemble the SGA-2025 parent sample through a sequence of discrete
processing stages, each building on the output of the one before:
merging six external source lists into a single, uncut candidate
catalog; removing contaminants and correcting errors through iterative
visual inspection; vetoing spurious detections with a
self-supervised-learning classifier and applying a battery of
additional photometric and positional cuts; masking the footprint
around bright Gaia and Tycho stars; and, finally, assembling the
released parent sample, release by release, from a versioned sequence
of hand-curated corrections.  This multi-catalog strategy responds
directly to a lesson learned building the SGA-2020, whose parent sample
HyperLeda \citep{makarov14a} supplied almost in its entirety.  A
completeness check against the independent HECATE catalog
\citep{kovlakas21a} found the SGA-2020 highly complete in practice
($\gtrsim\!95\%$ for galaxies larger than $\sim\!25\arcsec$), but traced
its residual, genuine misses and the pronounced surface-density
variations across its footprint to heterogeneity and query artifacts
inherited from relying on HyperLeda alone \citep{moustakas23a}.  Because
the mechanics of these
intermediate stages are logistically detailed but not, in themselves,
central to interpreting the resulting catalog, we summarize each briefly
below and relegate the full accounting -- every hand-curated resource we
draw on, every numerical threshold we apply, and every internal
consistency check we impose -- to Appendix~\ref{sec:appendix:parent}.

We build the no-cuts candidate catalog (Table~\ref{tab:parent-external})
by merging six external source lists, each contributing a distinct
piece of the final galaxy inventory.  HyperLeda \citep{makarov14a} is
our primary input, supplying $D_{25}$ diameters and positions for the
great majority of candidates; we reconcile it against NED-LVS
\citep[NED Local Volume Sample;][]{cook23a}, an independently assembled
NED cross-identification and redshift compilation, to establish a
high-confidence overlap sample and to recover HyperLeda objects that
NED resolves under a different name or position.  We supplement these
two all-sky compilations with the Local Volume Database \citep[LVD;][]
{pace25a}, a hand-vetted census of confirmed dwarf and Local Group
galaxies within $\sim\!20$~Mpc that supplies the most reliable
positions, distances, and structural parameters for this population;
every SGA-2020 galaxy \citep{moustakas23a} within the DR11 footprint,
carried forward directly; a hand-curated ``custom'' catalog of one-off
additions dominated by the SMUDGes ultra-diffuse and
low-surface-brightness galaxy sample \citep{zaritsky23a}; and a
supplemental list of large galaxies identified during earlier
(DR9/DR10) SGA processing.  We describe the cross-matching,
source-priority ordering, and geometry standardization we apply to
merge these six catalogs into a single, consistent system of
positions, diameters, axis ratios, and magnitudes in
Appendix~\ref{sec:input:parent:external}.

Beginning from this no-cuts catalog, we apply three further rounds of
cleanup before assembling the final parent sample: an iterative,
visual-inspection-driven pass that removes non-galaxy contaminants and
corrects erroneous positions, geometries, and cross-identifications
(Appendix~\ref{sec:parent:vicuts}); a self-supervised-learning
classifier that vetoes spurious detections, together with additional
close-pair resolution, Local Group/Galactic-contaminant flagging, and a
photometric cut on the smallest candidates (Appendix~\ref{sec:parent:archive});
and the construction of a Gaia+Tycho bright-star mask, which we use to
flag objects near or embedded within a bright star's masking radius for
special handling during ellipse-fitting (Appendix~\ref{sec:parent:gaiamask}).
Together these stages reduce the no-cuts catalog from $\sim\!5.10$~million
candidates to the $\sim\!1.19$~million-object archive-stage catalog we
carry forward into the final assembly below.  \todo{Update these object
  counts once the final (v1.0) processing version is available.}

We assemble the final parent sample from this archive-stage catalog not
in a single pass, but across a sequence of thirteen numbered releases
(v0.30 through v1.6, Table~\ref{tab:parent-overlays}), each of which
refines the previous release's fitted ellipse geometry via an
``overlay'' model: every release is defined entirely by a small,
versioned set of hand-curated CSV files recording the objects added,
dropped, individually corrected, or re-flagged at that stage of visual
inspection, so that the full history of the catalog's evolution
survives as a permanent, inspectable record rather than being folded
irreversibly into the catalog itself.  After applying each release's
overlay files, we re-run group-finding on the full, corrected catalog to
identify galaxy groups and assign each a primary member
(Section~\ref{sec:parent:groups}), and re-populate the \texttt{SAMPLE}
and \texttt{ELLIPSEMODE} bits (Tables~\ref{tab:sample-bitmask}
and~\ref{tab:ellipsemode-bitmask}) that downstream ellipse-fitting and
cataloging depend on.  In total, we made more than 45,000 individual
visual-inspection-driven additions, removals, or corrections across the
v0.30--v1.6 releases; \texttt{v1.6} is the final parent-sample version
adopted for the SGA-2025 release.  We describe this overlay bookkeeping,
and the release-specific corrections it required, in full in
Appendix~\ref{sec:parent:final}.

\subsection{Angular Group Catalog}
\label{sec:parent:groups}

We identify galaxy groups -- physically associated or angularly
overlapping systems whose members we model jointly in a single mosaic
(Section~\ref{sec:pipeline:ellipse}) -- by re-running a dedicated
group-finding algorithm at every parent-sample release
(Section~\ref{sec:parent:summary}), assigning each object a
\texttt{GROUP\_NAME}, a multiplicity (\texttt{GROUP\_MULT}), a
diameter-weighted center (\texttt{GROUP\_RA}/\texttt{GROUP\_DEC}/
\texttt{GROUP\_DIAMETER}), and a Boolean \texttt{GROUP\_PRIMARY} flag
marking the dominant (largest-diameter) member that most SGA-2025
analyses key on.  We treat two populations differently: objects flagged
\texttt{RESOLVED} (Milky Way satellites effectively resolved into
individual stars) or \texttt{FORCEPSF} (forced to a point-source
Tractor model) are excluded from automated linking entirely and instead
assigned trivially to their own singleton, one-member group
(\texttt{SGA.groups.make\_singleton\_group}), since neither
population's angular extent is physically meaningful for grouping
purposes; every other object is passed to the full linking algorithm
(\texttt{SGA.groups.build\_group\_catalog}).

\texttt{build\_group\_catalog} links objects using a threshold set to
the diameter-dependent scale of the pair being tested, rather than a
single fixed angle.  We first bucket the input catalog into coarse
spherical preclusters via \texttt{pydl}'s \texttt{spheregroup}
algorithm, using a generous $5\arcmin$ link length that cheaply groups
together any objects that could plausibly belong to the same physical
system while keeping the subsequent, more expensive pairwise test
confined to small preclusters.  Within each multi-member precluster we
then test every candidate pair against a size-dependent linking
threshold: for two ``small'' objects (diameter $<3\arcmin$) we compare
their center-to-center separation against $\sim\!1.6\times$ the mean of
their \emph{directional} radii -- the radius each ellipse presents
toward its neighbor, so that two elongated galaxies aligned end-to-end
link at a larger separation than two aligned side-by-side; for two
``big'' objects we instead use their circular (position-angle-independent)
radii scaled by $1.3\times$, avoiding spuriously large directional radii
for the most elongated systems; and for a big-small pair we use
directional radii scaled by $1.5\times$, floored at $40\%$ of the small
object's circular radius so that a small companion close to a large
host's isophotal edge is not missed.  Independent of this
threshold, we also link any pair whose ellipses geometrically overlap
once expanded by a $75\%$ margin, which catches close pairs the
distance-based test alone would miss; we floor every threshold at
$36\arcsec$ and reject linking altogether beyond $5\arcmin$ separation.

Because the coarse preclustering and pairwise linking above can still
leave near-duplicate or narrowly-missed groups unmerged, we apply two
further passes before finalizing the catalog.  First, we merge any pair
of provisional groups whose members' \emph{unscaled} ellipses overlap
but whose separation exceeded the $5\arcmin$ linking cutoff -- the case,
for example, of a small companion projected deep within a much larger
host's disk, too close in projection for the coarse preclustering step
to have placed them together but unambiguous once ellipse overlap is
checked directly.  Second, we compute each provisional group's
diameter-weighted center and merge any two group centers that land
within $52\arcsec$ of one another, catching residual duplicate groups
the ellipse-overlap check does not resolve; we also support a small,
hand-curated list of named object pairs to force-merge regardless of
geometry, for the rare cases visual inspection identifies as physically
associated but geometrically disjoint.  Finally, for each resulting
group we compute its adopted center as the diameter-weighted Cartesian
mean of its members' positions, its diameter as the maximum member
extent (separation from the group center plus half that member's own
diameter), measured from that center and floored at $30\arcsec$, and
its primary flag as the single largest-diameter member.

We invoke group-finding after every parent-sample release's overlay
corrections are applied (Section~\ref{sec:parent:summary}), so that the
additions, drops, and geometry updates at each release can, in
principle, reshape the group catalog; we then harmonize each group's
imaging-region bits and set the \texttt{OVERLAP} \texttt{SAMPLE} bit for
every member whose ellipse overlaps that of another member of its own
group (Appendix~\ref{sec:parent:final}).

%%----------------------------------------------------------------------
\section{Photometric Pipeline}
\label{sec:pipeline}
%%----------------------------------------------------------------------

\subsection{Imaging Data}
\label{sec:imaging}

The DESI Legacy Imaging Surveys \citep{dey19a} combine three public
imaging programs -- the DECam Legacy Survey (DECaLS), the
Beijing--Arizona Sky Survey (BASS), and the Mayall $z$-band Legacy
Survey (MzLS) -- to cover the extragalactic sky in $g$, $r$, $i$, and
$z$ over \texttt{dr11-south} and in $g$, $r$, and $z$ over
\texttt{dr11-north}.  The Dark Energy Camera \citep[DECam;][]{flaugher15a},
a 3~deg$^2$ mosaic imager with a pixel scale of $0\farcs262$ on the
Blanco 4-m telescope at CTIO, provides the DECaLS imaging; by DR11, the
\texttt{dr11-south} footprint aggregates not only DECaLS's own exposures
but also public archival DECam data from the Dark Energy Survey, the
DECam Local Volume Exploration Survey, and the DECam eROSITA Survey, all
reprocessed with a common pipeline to achieve uniform depth and
photometric calibration over the combined footprint.  We image
\texttt{dr11-north} by combining $g$- and $r$-band imaging from BASS
(the 90Prime imager on the Steward Observatory 2.3-m Bok telescope;
\citealt{zou17a}) with $z$-band imaging from MzLS (the Mosaic-3 imager
on the KPNO 4-m Mayall telescope).

We detect and model sources across all three imaging programs jointly
with the Tractor forward-modeling pipeline (Section~\ref{sec:pipeline:tractor}),
which poses source characterization as a $\chi^2$ minimization applied
directly to the pixel data from every overlapping exposure -- extracting
the full depth and resolution of the combined dataset without resorting
to heuristic deblending -- using spatially varying point-spread-function
models built with PSFEx and astrometric and photometric zero-points tied,
CCD by CCD, to Gaia and Pan-STARRS1.  DR11 is the most complete and
uniform Legacy Surveys release to date, covering $\sim\!31{,}000$~deg$^2$
across both hemispheres to typical $5\sigma$ point-source depths of
$g\approx24.7$, $r\approx24.2$, $i\approx23.8$, and $z\approx23.3$~mag.
\todo{Confirm final DR11 depths and area against the published DR11
  paper once available; the draft we consulted still marks these
  values TBC.}  Notably, DR11 is also the first Legacy Surveys release
to incorporate the SGA-2025 itself: our galaxy geometries are pre-loaded
as fitting priors (\texttt{REF\_CAT} $=$ \texttt{`L4'}) so that Tractor
treats each large galaxy as a single system rather than fragmenting its
disk into spurious compact sources or absorbing its outer envelope into
the sky background, a reciprocal dependence we revisit in
Section~\ref{sec:pipeline}.

%\subsection{Mid-Infrared: unWISE}
%\label{sec:imaging:wise}

We incorporate mid-infrared imaging from unWISE \citep{lang14a,
  meisner21a}, full-depth coadds of the Wide-field Infrared Survey
Explorer \citep[WISE;][]{wright10a} in bands $W1$ (3.4~$\mu$m), $W2$
(4.6~$\mu$m), $W3$ (12~$\mu$m), and $W4$ (22~$\mu$m).  The unWISE
coadds we use incorporate every NEOWISE epoch through the end of the
mission, providing the deepest possible $W1$/$W2$ coverage, while $W3$
and $W4$ retain the original full-mission depth from the cryogenic
phase of the WISE mission.  The unWISE coadds provide nearly uniform
coverage across the sky at angular resolutions of
$6\arcsec$--$12\arcsec$, set by the WISE point-spread function.

%\subsection{Ultraviolet: GALEX}
%\label{sec:imaging:galex}

We obtain ultraviolet photometry from the Galaxy Evolution Explorer
\citep[GALEX;][]{martin05a} GR6/GR7 data release in the FUV (1528~\AA{})
and NUV (2271~\AA{}) bands.  Unlike the optical and mid-infrared imaging,
GALEX coverage is non-uniform and does not extend over the full DR11
footprint; approximately $X$\% of SGA-2025 galaxies have at least one
GALEX band available.

\todo{Is it worth building an optical, GALEX, and unWISE depth map figure?}

\subsection{Image Coadds and Tractor Photometry}
\label{sec:pipeline:tractor}

For each galaxy group we build a custom set of image coadds and source
catalogs by invoking \textsl{legacypipe}'s \texttt{runbrick} pipeline
(\texttt{SGA.coadds.custom\_coadds}) on a synthetic ``brick'' centered on
the group position, rather than processing the group piecemeal across
whichever standard DR11 bricks happen to overlap it.  We size each
mosaic adaptively from the group's diameter, multiplicity, and axis
ratio -- inflating the width for small, elongated, or high-multiplicity
systems so that the full light profile and an adequate sky annulus are
always captured -- at the survey's native optical pixel scale
($0\farcs262$~pixel$^{-1}$), together with matched unWISE
($2\farcs75$~pixel$^{-1}$) and, where available, GALEX
($1\farcs5$~pixel$^{-1}$) coadds built in the same run.  \textsl{legacypipe}
and Tractor \citep{lang16a} are described in detail elsewhere
\citep{dey19a,moustakas23a}; we summarize here only the processing
options we adopt that depart from a standard DR11 brick run.

We fit the great majority of groups on their coadded images rather than
on the individual per-exposure frames (\texttt{--fit-on-coadds
--no-ivar-reweighting}), with iterative source-detection refinement
disabled (\texttt{--no-iterative}) in favor of a single detection pass at
fixed significance and saddle-point thresholds ($n\sigma=6$, saddle
fraction $=0.2$, saddle minimum $=4$).  We also depart from
\textsl{legacypipe}'s default handling of previously cataloged large
galaxies: ordinarily, sources falling within the aperture of a known
large galaxy (\texttt{REF\_CAT}~$=$~\texttt{`L4'}, i.e., the SGA itself;
Section~\ref{sec:imaging}) are forced to a point-source model to avoid
spurious multi-component fits, but since our custom mosaics exist
specifically to fit those same large galaxies in detail, we disable this
suppression by default (\texttt{--no-galaxy-forcepsf}); we leave it
enabled only for the small number of objects flagged \texttt{FORCEPSF}
in \texttt{ELLIPSEMODE} (Table~\ref{tab:ellipsemode-bitmask}), which
forces every source in that mosaic to a point-source model.  Finally, we
mask bright sources that fall outside the blob being fit
(\texttt{--bright-masking}) and apply custom segmentation around large
galaxies (\texttt{--galaxy-masking}), to keep neighboring bright stars
and galaxies from contaminating detection and model-fitting.

Two categories of object bypass this pipeline entirely.  Objects flagged
\texttt{RESOLVED} in \texttt{ELLIPSEMODE} -- Milky Way satellites so
nearby that they resolve into individual stars rather than a smooth
light profile -- skip Tractor altogether; we instead cut out existing
coadd imagery directly, at 100$\times$ the standard pixel scale for the
very largest such systems (the LMC, SMC, and Antlia~II) and 50$\times$
for the next largest (Crater~II and Bootes~III), since these objects
span degrees on the sky and a native-resolution mosaic would be both
unnecessary and computationally prohibitive.  Objects for which we
request \texttt{--skip-tractor} are likewise reduced to simple image
cutouts, retaining an inverse-variance map so that photometry can still,
in principle, be attempted later; we flag these objects
\texttt{SKIPTRACTOR} in \texttt{ELLIPSEBIT}
(Table~\ref{tab:ellipsebit-bitmask}).

\subsection{Ellipse Fitting}
\label{sec:pipeline:ellipse}

\JM{Stub refs, placement TBD: Table~\ref{tab:sample-bitmask} (\texttt{SAMPLE}),
Table~\ref{tab:ellipsemode-bitmask} (\texttt{ELLIPSEMODE}),
Table~\ref{tab:ellipsebit-bitmask} (\texttt{ELLIPSEBIT}),
Table~\ref{tab:maskbits-common} and Table~\ref{tab:maskbits-bands}
(\texttt{OPTMASKBITS}/\texttt{GALEXMASKBITS}/\texttt{UNWISEMASKBITS}).}
\begin{deluxetable*}{lcl}
\tablecaption{\texttt{SAMPLE} Bitmask Definitions
  \label{tab:sample-bitmask}}
\tablewidth{\textwidth}
\tablehead{\colhead{Bit Name} & \colhead{Value} & \colhead{Description}}
\startdata
\texttt{LVD}      & 1  & \parbox{4.0in}{Source from the Local Volume Database of nearby galaxies (Section~\ref{sec:parent}).} \\
\texttt{MCLOUDS}  & 2  & \parbox{4.0in}{Located within the Magellanic Cloud footprint.} \\
\texttt{GCLPNE}   & 4  & \parbox{4.0in}{Located within a globular cluster or planetary nebula mask.} \\
\texttt{NEARSTAR} & 8  & \parbox{4.0in}{Near a bright star (star--galaxy separation $<1.2\times$ the Gaia/Tycho mask radius).} \\
\texttt{INSTAR}   & 16 & \parbox{4.0in}{Inside a bright star (star--galaxy separation $<0.5\times$ the Gaia/Tycho mask radius).} \\
\texttt{OVERLAP}  & 32 & \parbox{4.0in}{Initial ellipse overlaps another (initial) SGA-2025 ellipse.} \\
\enddata
\tablecomments{Set in the parent catalog (Section~\ref{sec:parent}) and propagated to the merged catalog; a source may have more than one bit set simultaneously.}
\end{deluxetable*}

\begin{deluxetable*}{lcl}
\tablecaption{\texttt{ELLIPSEMODE} Bitmask Definitions
  \label{tab:ellipsemode-bitmask}}
\tablewidth{\textwidth}
\tablehead{\colhead{Bit Name} & \colhead{Value} & \colhead{Description}}
\startdata
\texttt{FIXGEO}       & 1   & \parbox{4.0in}{Fix the ellipse geometry to its initial values (no free fitting).} \\
\texttt{RESOLVED}     & 2   & \parbox{4.0in}{Milky Way satellite, effectively resolved into individual stars; Tractor fitting is not performed.} \\
\texttt{FORCEPSF}     & 4   & \parbox{4.0in}{Force all Tractor sources in the mosaic to be fitted as point sources.} \\
\texttt{FORCEGAIA}    & 8   & \parbox{4.0in}{Force Tractor to fit only Gaia sources (not used in the SGA-2025 processing).} \\
\texttt{LESSMASKING}  & 16  & \parbox{4.0in}{Apply less aggressive neighbor masking during ellipse-fitting (Section~\ref{sec:pipeline:ellipse:mask}).} \\
\texttt{MOREMASKING}  & 32  & \parbox{4.0in}{Apply more aggressive neighbor masking during ellipse-fitting.} \\
\texttt{MOMENTPOS}    & 64  & \parbox{4.0in}{Adopt the moment-based center instead of the Tractor position.} \\
\texttt{TRACTORGEO}   & 128 & \parbox{4.0in}{Adopt the Tractor geometry instead of the ellipse-fit geometry.} \\
\texttt{NORADWEIGHT}  & 256 & \parbox{4.0in}{Disable radial weighting when determining the light-weighted geometry.} \\
\enddata
\tablecomments{Requested prior to processing; describes the intended fitting strategy for a given object (Section~\ref{sec:pipeline:ellipse}). Populated in part from the hand-curated resources of Table~\ref{tab:parent-archive} and consolidated across releases in each overlay's \texttt{flags.csv} (Section~\ref{sec:parent:final}).}
\end{deluxetable*}

\begin{deluxetable*}{lcl}
\tablecaption{\texttt{ELLIPSEBIT} Bitmask Definitions
  \label{tab:ellipsebit-bitmask}}
\tablewidth{\textwidth}
\tablehead{\colhead{Bit Name} & \colhead{Value} & \colhead{Description}}
\startdata
\texttt{NOTRACTOR}          & 1     & \parbox{4.0in}{No Tractor model available (object dropped or never fitted).} \\
\texttt{TRACTORPSF}         & 2     & \parbox{4.0in}{Tractor fit this source as a point source (PSF model).} \\
\texttt{FIXGEO}             & 4     & \parbox{4.0in}{The ellipse geometry was held fixed rather than freely fitted.} \\
\texttt{BLENDED}            & 8     & \parbox{4.0in}{The source center falls within the ellipse of another SGA-2025 source.} \\
\texttt{LARGESHIFT}         & 16    & \parbox{4.0in}{Large positional shift between the initial and moment-refined ellipse center.} \\
\texttt{LARGESHIFT\_TRACTOR}& 32    & \parbox{4.0in}{Large positional shift between the initial center and the Tractor position.} \\
\texttt{MAJORGAL}           & 64    & \parbox{4.0in}{Major galaxy requiring special processing treatment.} \\
\texttt{OVERLAP}            & 128   & \parbox{4.0in}{The fitted ellipse overlaps another SGA-2025 ellipse.} \\
\texttt{SATELLITE}          & 256   & \parbox{4.0in}{Satellite galaxy located within the ellipse of a larger system.} \\
\texttt{MOMENTPOS}          & 512   & \parbox{4.0in}{The moment-based center was adopted.} \\
\texttt{TRACTORGEO}         & 1024  & \parbox{4.0in}{The Tractor geometry was adopted.} \\
\texttt{NORADWEIGHT}        & 2048  & \parbox{4.0in}{Radial weighting was not applied during isophote fitting.} \\
\texttt{LESSMASKING}        & 4096  & \parbox{4.0in}{Less aggressive neighbor masking was applied.} \\
\texttt{MOREMASKING}        & 8192  & \parbox{4.0in}{More aggressive neighbor masking was applied.} \\
\texttt{FAILGEO}            & 16384 & \parbox{4.0in}{Isophote fitting failed; \texttt{D26} may be unreliable.} \\
\texttt{SKIPTRACTOR}        & 32768 & \parbox{4.0in}{Tractor fitting was skipped entirely for this mosaic.} \\
\enddata
\tablecomments{Set after processing; records the fitting outcome and data quality actually realized for a given object (Section~\ref{sec:pipeline:ellipse}), as opposed to \texttt{ELLIPSEMODE} (Table~\ref{tab:ellipsemode-bitmask}), which records the intended strategy requested beforehand.}
\end{deluxetable*}

\begin{deluxetable*}{lcl}
\tablecaption{Shared Low-Order Bits of \texttt{OPTMASKBITS}, \texttt{GALEXMASKBITS}, and \texttt{UNWISEMASKBITS}
  \label{tab:maskbits-common}}
\tablewidth{\textwidth}
\tablehead{\colhead{Bit Name} & \colhead{Value} & \colhead{Description}}
\startdata
\texttt{brightstar} & 1 & \parbox{4.0in}{Bright star, medium star, or star cluster (\textsl{legacypipe} \texttt{MASKBITS}).} \\
\texttt{gaiastar}   & 2 & \parbox{4.0in}{Gaia point source (\texttt{TYPE = PSF}).} \\
\texttt{galaxy}     & 4 & \parbox{4.0in}{Extended (non-reference) source.} \\
\texttt{reference}  & 8 & \parbox{4.0in}{SGA-2025 reference source.} \\
\enddata
\tablecomments{Per-band imaging mask bits stored in the per-group \texttt{-maskbits.fits[.fz]} files and the per-galaxy \texttt{-ellipse-\{dataset\}.fits} files (Section~\ref{sec:pipeline}). All three masks share this low-order encoding; the higher-order, per-band bits are listed separately in Table~\ref{tab:maskbits-bands}.}
\end{deluxetable*}

\begin{deluxetable}{lccc}
\tablecaption{Per-Band Bits of \texttt{OPTMASKBITS}, \texttt{GALEXMASKBITS}, and \texttt{UNWISEMASKBITS}
  \label{tab:maskbits-bands}}
\tablehead{\colhead{Value} & \colhead{\texttt{OPTMASKBITS}} & \colhead{\texttt{GALEXMASKBITS}} & \colhead{\texttt{UNWISEMASKBITS}}}
\startdata
16  & $g$   & FUV      & W1 \\
32  & $r$   & NUV      & W2 \\
64  & $i$   & \nodata  & W3 \\
128 & $z$   & \nodata  & W4 \\
\enddata
\tablecomments{Each mask assigns its own bands to bits 16, 32, 64, and 128, on top of the shared low-order bits of Table~\ref{tab:maskbits-common}; \textsl{griz} imaging uses all four band bits, while GALEX (FUV/NUV) uses only the first two.}
\end{deluxetable}


We perform ellipse photometry on entire galaxy groups at once, rather
than on each galaxy in isolation, using a three-stage iterative
procedure closely related to the methodology we describe in
\citet{moustakas23a}: we first assemble each group's imaging and source
catalogs (Section~\ref{sec:pipeline:ellipse:read}); we then iteratively
refine every member's elliptical geometry and construct a hierarchy of
source masks (Section~\ref{sec:pipeline:ellipse:mask}); and we finally
fit the resulting surface-brightness profiles in two passes across all
available imaging datasets (Section~\ref{sec:pipeline:ellipse:fit}).  We
adopt a group-level architecture because SGA galaxies frequently overlap
physically associated or line-of-sight companions -- fitting each galaxy
independently, blind to its neighbors' light, would systematically bias
the recovered geometry and isophotal radii; instead, every stage below
operates simultaneously on all group members, so that the mask we build
for one object accounts for the position, geometry, and flux of every
other cataloged source in the same mosaic.

\subsubsection{Imaging Preparation}
\label{sec:pipeline:ellipse:read}

For each group, we first read the coadded optical, unWISE, and (where
available) GALEX mosaics -- images, inverse-variance maps, PSF models,
and Tractor model images -- together with the full Tractor catalog for
the mosaic footprint and the group's sample catalog of SGA targets.  We
construct a bright-star mask from the \texttt{MASKBITS} image (the
\texttt{BRIGHT}, \texttt{MEDIUM}, and \texttt{CLUSTER} bits) together
with an independent Gaia/Tycho-based core mask, and we match each SGA
object in the sample catalog to its corresponding Tractor source by
reference ID.  Galaxies large or diffuse enough to require visual
inspection bypass this matching, and every fitting stage that follows,
entirely; we instead populate their geometry directly from the parent
catalog.

\subsubsection{Masking and Geometry Refinement}
\label{sec:pipeline:ellipse:mask}

Beginning from each object's initial position and geometry -- taken from
the parent catalog or, where preferred, from the Tractor model -- we
iteratively refine the elliptical geometry from light-weighted second
moments of the image, alternating between remeasuring the moments and
rebuilding the elliptical mask used to compute them until the solution
converges (Table~\ref{tab:ellipsefit-params}).  We weight the moments
radially, down-weighting the galaxy's core relative to its outskirts, so
that the recovered semi-major axis, axis ratio, and position angle trace
the extended, low-surface-brightness envelope rather than the often
rounder, bulge-dominated central light.  We flag objects whose
moment-based center shifts by more than $\Delta_{\rm max}$ from its
starting position as having an unreliable geometry solution, and we flag
objects whose semi-major axis falls below $f_{\rm sat}$ times that of
the largest overlapping neighbor as physically blended satellites
requiring closer scrutiny.

Once the geometry converges, we build a final mask for each object by
combining, in addition to the per-band data-quality mask: (1) the
bright-star mask described above; (2) the elliptical footprint of every
other SGA source in the group, protected out to $f_{\rm ref}$ times its
own semi-major axis; and (3) the footprints of neighboring
Tractor-detected galaxies, which we classify as major companions --
masked unconditionally, including within the target's own ellipse -- if
their flux exceeds $F_{\rm major}^{\rm final}$ times the target's, and
otherwise, optionally, as minor companions masked only outside the
target's own ellipse.  We then subtract the Tractor model of every
masked source, so that the images entering the fitting stage
(Section~\ref{sec:pipeline:ellipse:fit}) contain only the flux of the
target galaxy and any unmasked, unmodeled structure; we propagate the
same reference/bright-star/companion masking, rescaled to the
appropriate pixel scale, to the unWISE and GALEX mosaics, so that the
optical, mid-infrared, and ultraviolet photometry we describe in
Section~\ref{sec:pipeline:cog} share a consistent masking footprint.

\todo{Figure: illustrate the mask hierarchy for one representative
  group, e.g., three panels showing (a) the initial, Tractor-based
  mask, (b) the moment-refined mask after geometry convergence, and (c)
  the final mask, with major/minor companions and bright stars masked
  and their Tractor models subtracted.}

Table~\ref{tab:ellipsefit-params} summarizes the free parameters
governing the masking and fitting stages described in this section,
together with their adopted default values.

\begin{deluxetable*}{lcl}
\tablecaption{Free Parameters of the Ellipse-Fitting Pipeline
  \label{tab:ellipsefit-params}}
\tablewidth{\textwidth}
\tablehead{\colhead{Parameter} & \colhead{Default} & \colhead{Description}}
\startdata
$n_{\rm geo}$               & 2                  & \parbox{4.0in}{Number of moment-based geometry-refinement iterations per object.} \\
$F_{\rm major}^{\rm geo}$   & 0.01               & \parbox{4.0in}{Minimum flux ratio (relative to the target) for a neighbor to be classified ``major'' during geometry refinement.} \\
$F_{\rm major}^{\rm final}$ & 0.1                & \parbox{4.0in}{As above, applied when building the final mask after the geometry has converged.} \\
$f_{\rm ref}$               & 2.0                & \parbox{4.0in}{Factor by which another SGA source's semi-major axis is extended when protecting its flux from masking.} \\
$\alpha_{\rm rad}$          & 0.7                & \parbox{4.0in}{Exponent of the radial weight applied to the flux in the second-moment calculation.} \\
$f_{\rm sat}$               & 0.3                & \parbox{4.0in}{Semi-major-axis ratio (relative to the largest overlapping neighbor) below which an object is flagged a satellite.} \\
$\Delta_{\rm max}$          & 3\farcs5           & \parbox{4.0in}{Maximum allowed shift between the initial and moment-refined center before flagging an unreliable solution.} \\
$\mu_{\rm iso}$             & 22, 23, 24, 25, 26 & \parbox{4.0in}{Surface-brightness thresholds (mag~arcsec$^{-2}$) for isophotal-radius measurements.} \\
$N_{\rm MC}$                & 50                 & \parbox{4.0in}{Number of Monte Carlo realizations for isophotal-radius uncertainties.} \\
\enddata
\end{deluxetable*}


\subsubsection{Two-Pass Isophotal Fitting}
\label{sec:pipeline:ellipse:fit}

We fit the masked, model-subtracted mosaics in two passes.  In the first
pass, we build an aggressive initial mask -- classifying companions as
major down to a flux ratio of just $F_{\rm major}^{\rm geo}$ and masking
minor companions throughout -- and we fit only the optical bands,
without Monte Carlo resampling, to obtain a first-pass estimate of the
$\mu = 26$~mag~arcsec$^{-2}$ isophotal radius, $R(26)$, for every object
in the group.  In the second pass, we rebuild the mask using this
refined geometry -- relaxing the major-companion threshold to
$F_{\rm major}^{\rm final}$ and, by default, leaving minor companions
unmasked -- and we fit the full complement of requested imaging datasets
(optical, unWISE, and GALEX) simultaneously.

For each object, band, and dataset, we measure elliptical-aperture
surface photometry on a grid of semi-major axes using the
isophote-fitting machinery of \texttt{photutils} \citep{bradley23a},
sigma-clipping outlier pixels within each isophote before averaging; we
anchor the sampling grid to the object's moment-refined semi-major axis
and rescale it, band to band and dataset to dataset, to hold a roughly
constant number of pixels -- and hence roughly constant signal-to-noise
-- per annulus.  We derive isophotal radii at surface-brightness levels
of $\mu = 22$, 23, 24, 25, and 26~mag~arcsec$^{-2}$ in each available
optical band from a Monte Carlo resampling of the surface-brightness
profile ($N_{\rm MC}$ realizations by default), which we use to assign
an asymmetric uncertainty from the interquartile range of the resulting
radius distribution rather than propagating formal photometric errors
alone; we adopt the $\mu = 26$~mag~arcsec$^{-2}$ isophote in the $r$
band as the SGA-2025 diameter, $D_{26}$ (arcmin).  We derive the
curve-of-growth photometry and half-light radii that we report in all
ten photometric bands from these same per-annulus measurements, as we
describe next in Section~\ref{sec:pipeline:cog}.

\subsection{Curve-of-Growth Photometry}
\label{sec:pipeline:cog}

In addition to the discrete isophotal radii we derive in
Section~\ref{sec:pipeline:ellipse:fit}, we characterize each object's
full surface-brightness profile with two complementary flux
measurements, computed simultaneously with the isophotal fit and from
the identical per-annulus aperture photometry: a set of fixed reference
apertures, and a parametric fit to the curve of growth that yields an
asymptotic total magnitude and a half-light radius.  We measure both for
every object, band, and imaging dataset (optical, unWISE, and GALEX), so
that \texttt{COG\_MTOT\_$\langle$band$\rangle$} and
\texttt{SMA50\_$\langle$band$\rangle$} are defined consistently across
all ten photometric bands.

For each band, we first measure elliptical-aperture photometry at a
small set of fixed reference apertures, scaled to the object's
moment-refined semi-major axis by factors of 0.5, 1, 1.25, 1.5, 2, and
3; we record the resulting flux, flux uncertainty, and masked-pixel
fraction at each aperture as \texttt{FLUX\_AP$\langle$NN$\rangle$\_$\langle$band$\rangle$},
\texttt{FLUX\_ERR\_AP$\langle$NN$\rangle$\_$\langle$band$\rangle$}, and
\texttt{FMASKED\_AP$\langle$NN$\rangle$\_$\langle$band$\rangle$}, providing a
model-independent photometric anchor at a common physical scale across
the sample, independent of the parametric fit we describe next.

We complement these fixed apertures with a parametric fit to the full
curve of growth.  Using the same per-annulus aperture-flux measurements
underlying the isophotal fit (Section~\ref{sec:pipeline:ellipse:fit}),
we model the enclosed-flux profile, expressed as an AB magnitude $m(r)$
as a function of semi-major axis $r$, as
\begin{equation}
  m(r) = m_{\rm tot} + \Delta m \left[1 - \exp\!\big(-\alpha_1
    (r/r_0)^{-\alpha_2}\big)\right],
  \label{eq:cog}
\end{equation}
where $m_{\rm tot}$ is the asymptotic total magnitude as $r \to \infty$,
$\Delta m$ is the (positive) amplitude by which the innermost-aperture
magnitude exceeds $m_{\rm tot}$, $\alpha_1$ and $\alpha_2$ jointly set
the curve's normalization and steepness, and we fix the scale radius
$r_0$ to the object's moment-refined semi-major axis.  We fit
Equation~\ref{eq:cog} in terms of $(m_{\rm tot}, \Delta m, \ln\alpha_1,
\ln\alpha_2)$ rather than $(m_{\rm tot}, \Delta m, \alpha_1, \alpha_2)$
directly -- since $\alpha_1$ and $\alpha_2$ are defined only for
positive values, optimizing their logarithms instead lets the
nonlinear solver explore the full real line without encountering a
positivity boundary, which we find stabilizes convergence considerably
across the wide range of profile shapes spanned by the sample.  We
solve for these four parameters via a robust (soft-$L_1$-loss)
nonlinear least-squares minimization with an analytic Jacobian,
weighting each aperture by its flux-derived magnitude uncertainty --
with a 0.02~mag floor added in quadrature so that the innermost,
highest signal-to-noise apertures cannot dominate the fit -- and we
estimate parameter uncertainties from the fit's covariance matrix,
scaled by the reduced chi-square.  We report the asymptotic total
magnitude, $m_{\rm tot}$, as \texttt{COG\_MTOT\_$\langle$band$\rangle$}
in the AB system for all ten photometric bands: $G$, $R$, $I$, $Z$
(optical), $W1$--$W4$ (unWISE), and FUV, NUV (GALEX), together with its
formal uncertainty.

From the same fitted parameters we analytically invert
Equation~\ref{eq:cog} for the semi-major axis $r_f$ enclosing an
arbitrary fraction $f$ of the total flux -- i.e., where $m(r_f) -
m_{\rm tot} = -2.5\log_{10} f \equiv \delta_f$ -- to obtain
\begin{equation}
  r_f = r_0 \left[\frac{\alpha_1}{-\ln\!\left(1 - \delta_f/\Delta
      m\right)}\right]^{1/\alpha_2},
  \label{eq:rfrac}
\end{equation}
and we propagate the fit's covariance matrix through
Equation~\ref{eq:rfrac} via its analytic partial derivatives with
respect to $(\Delta m, \ln\alpha_1, \ln\alpha_2)$ to obtain the
corresponding $1\sigma$ uncertainty on $r_f$.  We evaluate
Equation~\ref{eq:rfrac} at $f = 0.5$ ($\delta_{0.5} \approx
0.753$~mag) to obtain the half-light radius, which we report as
\texttt{SMA50\_$\langle$band$\rangle$} together with its propagated
uncertainty; because $r_{0.5}$ follows directly from $(\Delta
m,\alpha_1,\alpha_2)$ rather than from an independent measurement, its
uncertainty reflects the full parameter covariance of the
curve-of-growth fit.

\todo{Figure: illustrate the curve-of-growth approach for one
  representative galaxy -- the aperture-flux (or surface-brightness)
  profile with the fitted model (Equation~\ref{eq:cog}) overlaid, the
  fixed reference apertures, the asymptotic $m_{\rm tot}$, and the
  resulting half-light radius.}

%%----------------------------------------------------------------------
\section{The SGA-2025 Catalog}
\label{sec:catalog}
%%----------------------------------------------------------------------

\subsection{Catalog Construction}
\label{sec:catalog:construction}

We assemble the released SGA-2025 catalog in two stages, separating the
purely internal products of our imaging pipeline (Section~\ref{sec:pipeline})
from the external cross-matching that turns those measurements into a
science-ready dataset.  In the first stage we merge, in parallel across MPI
ranks, the per-group ellipse-fit geometry, isophotal photometry, and
curve-of-growth measurements (Sections~\ref{sec:pipeline:ellipse} and
\ref{sec:pipeline:cog}) together with the row-matched forced Tractor
photometry into a single per-region ``beta'' catalog; at this stage every
object carries only the internal \texttt{SGAID} and the coordinates,
diameter, and name inherited from its parent-catalog entry
(Section~\ref{sec:parent})---no NED cross-identification, spectroscopic
redshift, or distance has yet been attached.  In the second stage we merge
the beta catalog with three external compilations: (1) NED positional and
by-name cross-matching, which supplies the adopted galaxy name
(\texttt{GALAXY}, \texttt{ALTNAMES}), morphological type, and NED's own
redshift and distance measurements; (2) DESI DR1 and SDSS DR17
spectroscopy, matched by position to every object within its fitted
ellipse; and (3) LVD velocities and distances for every Local Volume
Database object.  From these inputs we resolve a single adopted redshift
and distance per galaxy (Sections~\ref{sec:catalog:redshifts} and
\ref{sec:catalog:distances}) and train a photometric-redshift model
(Section~\ref{sec:catalog:redshifts}) on the resulting spectroscopic
subsample.  The output of this second stage is the public, value-added
\texttt{SGA2025-v1.0} catalog we describe in the remainder of this section.

\subsection{Catalog Structure}
\label{sec:catalog:structure}

We distribute the SGA-2025 as two FITS files, one per imaging region:

\begin{itemize}
  \item \texttt{SGA2025-dr11-south-v1.0.fits} ($\sim\!395{,}000$ entries)
  \item \texttt{SGA2025-dr11-north-v1.0.fits} ($\sim\!90{,}000$ entries)
\end{itemize}

Each file contains three binary table extensions.  The primary extension
(\texttt{SGA2025}) holds the science catalog described here: identification,
geometry, adopted redshifts and distances, and a nominal flux per band for
every galaxy.  The row-matched \texttt{ELLIPSEPHOT} extension carries the
full aperture-photometry and curve-of-growth grid underlying those nominal
fluxes (Section~\ref{sec:pipeline:cog}), keyed to \texttt{SGA2025} via
\texttt{SGAID}.  The row-matched \texttt{TRACTOR} extension contains
band-specific forced Tractor photometry quantities, keyed via
\texttt{REF\_ID} (duplicated as a leading \texttt{SGAID} column for
convenience).

\subsection{Key Columns}
\label{sec:catalog:columns}

Table~\ref{tab:columns} summarizes the most frequently used catalog
columns; we describe the adopted-redshift and adopted-distance columns
separately in Sections~\ref{sec:catalog:redshifts} and
\ref{sec:catalog:distances}.  We provide a complete data model in the
online documentation at \url{https://sga.readthedocs.io}.

% \input{tables/catalog_columns.tex}

\subsection{Data Access}
\label{sec:catalog:access}

We define the SGA-2025 sky coverage using two contiguous imaging regions
from the DR11 footprint (Figure~\ref{fig:sky}).  We image the southern
region (\texttt{dr11-south}; $\sim\!XX{,}000$~deg$^2$) with the Dark
Energy Camera (DECam) on the Cerro Tololo 4-m Blanco telescope in the
$g$, $r$, $i$, and $z$ bands; it contains $\sim\!363{,}000$ galaxies.
We image the northern region (\texttt{dr11-north};
$\sim\!XX{,}000$~deg$^2$; decl.\ $\gtrsim +32\arcdeg$) with the
Beijing--Arizona Sky Survey (BASS) and the Mayall $z$-band Legacy Survey
(MzLS) in $g$, $r$, and $z$; it contains $\sim\!82{,}000$ galaxies.
Together the two regions span $\sim\!30{,}000$~deg$^2$ and encompass
$\sim\!486{,}000$ resolved galaxies.  We describe the optical, mid-infrared,
and ultraviolet imaging that make up this dataset, in turn, in the three
subsections below.

\begin{figure*}
  \centering
  \includegraphics[width=\textwidth]{figures/sga2025-sky.png}
  \caption{Mollweide equal-area projection of the SGA-2025 galaxy surface
    density (galaxies deg$^{-2}$) in equatorial coordinates (J2000),
    centered at RA $= 120\arcdeg$.  The sample shown is restricted to
    group-primary objects ($N = 445{,}726$; one representative per galaxy
    group).  The solid black line traces the Galactic plane ($b = 0\arcdeg$)
    and the dashed lines mark $|b| = 10\arcdeg$, the approximate
    exclusion boundary applied during source selection.  The DR11 footprint
    defines two contiguous regions: the DECam south (low-latitude southern
    sky) and the BASS+MzLS north (decl.\ $\gtrsim +32\arcdeg$).}
  \label{fig:sky}
\end{figure*}

%\begin{figure*}
%  \centering
%  \includegraphics[width=\textwidth]{figures/sga2025-sky-south.png}
%  \includegraphics[width=\textwidth]{figures/sga2025-sky-north.png}
%  \caption{Galaxy surface density for the two imaging regions separately,
%    in the same projection as Figure~\ref{fig:sky}.
%    \textit{Top:} DR11 South imaged with DECam ($N = 363{,}633$).
%    \textit{Bottom:} DR11 North imaged with BASS+MzLS ($N = 82{,}093$).}
%  \label{fig:sky-regions}
%\end{figure*}

We make the catalog files, postage-stamp mosaics, and
ellipse-photometry outputs publicly available through the Legacy
Survey data portal at \url{https://sga.legacysurvey.org}.  We archive
the companion analysis code---which we use to generate the figures in
this paper---at Zenodo \todo{bib} and mirror it on GitHub at
\url{https://github.com/moustakas/sgapaper}.

\subsection{Validation}
\label{sec:catalog:validation}

We validate our ellipse-fit geometry against the WISE Extended Source
Catalog hundred sample \citep[WXSC-100;][]{jarrett19a}, an independent,
near-infrared reference sample of the 100 largest, brightest galaxies on
the sky.  We cross-match SGA-2025 against WXSC-100 by position and
compare, in Figure~\ref{fig:wxsc100}, our astrometry, isophotal diameters,
position angles, and ellipticities against the corresponding WXSC-100
measurements, color-coded by Hubble type.  We find good overall agreement
across all four quantities, with the residual scatter consistent with the
combination of measurement uncertainty and the intrinsically different
isophotal definitions ($\mu_{26}$ optical versus $W1$-band) underlying the
two catalogs.

\begin{figure*}
  \centering
  \includegraphics[width=\textwidth]{figures/sga2025-vs-wxsc100.png}
  \caption{Comparison of SGA-2025 geometry against the WXSC-100 sample
    \citep{jarrett19a}, color-coded by Hubble type.
    \textit{(a)} Astrometric offsets, $\Delta$R.A.\ and $\Delta$decl.
    \textit{(b)} Isophotal diameter $D(26)$ versus twice the WXSC-100
    $W1$-band radius, $2\,R_{W1}$.
    \textit{(c)} Position angle $\phi$.
    \textit{(d)} Ellipticity $\epsilon = 1-b/a$.
    In each panel the solid black line marks one-to-one agreement.}
  \label{fig:wxsc100}
\end{figure*}

\subsection{Redshifts}
\label{sec:catalog:redshifts}

We compile the adopted spectroscopic redshift (\texttt{Z}) from four
independent sources, in order of decreasing priority: LVD, DESI DR1, SDSS
DR17, and NED.  We adopt LVD \citep{pace25a} redshifts first, since the
LVD is a hand-curated compilation of velocities and distances for galaxies
within $\sim\!20$~Mpc and is the most reliable source available for the
nearby population it covers.  Failing an LVD measurement, we cross-match
against the DESI DR1 spectroscopic catalog \citep[DESI Collaboration et
al.~2025;][]{desi-collaboration26a}, restricting to spectra within
3.5\arcsec\ of the fitted center, discarding stellar and implausibly
high-redshift ($z > 1$) classifications as always spurious for this sample,
and preferring the highest-confidence galaxy spectrum with the largest
$\Delta\chi^2$ when several fall within the ellipse; we obtain redshifts
for $\sim\!X{,}000$ galaxies this way, predominantly through the DESI
Bright Galaxy Survey, and record the number of DESI spectra examined per
galaxy in \texttt{NSPEC\_DESI}.  We next cross-match against SDSS DR17
under comparable positional and quality cuts, independently of---and in
preference to---NED's own aggregated copy of SDSS spectroscopy, whose
quality flags NED does not propagate.  For the remainder we adopt NED's
redshift, replacing its aggregated ``fiducial'' value with one drawn from
its full redshift-measurement list whenever that value is suspiciously
high, missing, or drawn from a source we consider unreliable---most often
because a background source's redshift has been misattributed to a
foreground dwarf via fiber contamination---adopting the measurement from
whichever cluster of mutually consistent values dominates and flagging the
residual, ambiguous cases for visual inspection.

We store the adopted redshift, its inverse variance, and its source
(\texttt{Z}, \texttt{Z\_IVAR}, \texttt{Z\_REF} $\in \{$\texttt{LVD},
\texttt{DESI}, \texttt{SDSS}, \texttt{NED}$\}$); \texttt{Z\_IVAR} $= 0$
identifies a missing measurement, since \texttt{Z} $= 0$ is itself a valid,
if rare, redshift for the nearest galaxies.  Because DESI otherwise takes
priority whenever measured, we additionally override its adopted value in
two circumstances that would otherwise let an isolated discrepant spectrum
stand: when SDSS and NED are both measured and agree with each other but
not with DESI, and when an adopted DESI or SDSS redshift exceeds $z = 0.3$
---implausibly high for this overwhelmingly low-redshift sample---while
an independent NED measurement is on file.  We resolve the small residual
of flagged redshifts by hand and feed the results back into the catalog; we
provide the full override and quality-flag bookkeeping in the online
documentation.

Beyond this adopted redshift we derive two further quantities for every
galaxy.  We train a random-forest regressor \citep[\texttt{scikit-learn};]
[]{pedregosa11a} on the spectroscopic subsample, using extinction-corrected
optical/infrared colors, apparent magnitude, surface brightness, angular
diameter, axis ratio, and Tractor S\'ersic index as features, validate it
via 5-fold cross-validation, and predict a photometric redshift
(\texttt{Z\_PHOT}) for every galaxy, with its uncertainty
(\texttt{Z\_PHOT\_IVAR}) set from the spread of predictions across the
forest.  We also convert every measured redshift to the CMB frame and,
below 20,000~km~s$^{-1}$, remove the mean peculiar velocity predicted by
\texttt{pvhub} \citep{said20a}, propagating the model-to-model scatter into
the resulting \texttt{Z\_COSMO\_IVAR}; where no spectroscopic redshift
exists we substitute, in turn, the Hubble-flow redshift implied by an
independently measured distance (Section~\ref{sec:catalog:distances}) or,
failing that, the CMB-converted photometric redshift, so that
\texttt{Z\_COSMO} is defined as broadly across the sample as our ancillary
distance and photometric information allow.

We quantify the resulting redshift completeness in
Figure~\ref{fig:redshift-completeness}, where we show, as a function of
isophotal diameter $D(26)$, the fraction of all group members (not just
group primaries) with an adopted spectroscopic redshift, broken down by
source catalog.  We find that completeness rises steeply with increasing
angular size, as expected given the target selection of the legacy
spectroscopic surveys underlying NED and the DESI Bright Galaxy Survey;
at the largest diameters the catalog is nearly complete, while toward the
angular-size limit of the sample the NED and DESI contributions
increasingly dominate the incremental gain over LVD, whose reach is
confined to the Local Volume.  SDSS DR17 contributes a comparatively small
number of adopted redshifts, concentrated where DESI has not yet observed;
we fold it into the total completeness curve without a dedicated
breakdown.

\begin{figure}
  \centering
  \includegraphics[width=\columnwidth]{figures/redshift-completeness.png}
  \caption{Fraction of SGA-2025 galaxies (all group members, not just
    group primaries) with an adopted spectroscopic redshift as a function
    of isophotal diameter $D(26)$, broken down by source catalog: LVD
    (green), DESI DR1 (blue), and NED (orange).  The black curve shows
    the total completeness fraction, \texttt{Z\_IVAR} $> 0$, summed over
    all four sources, including the small SDSS DR17 contribution not
    broken out separately.  Completeness increases toward larger angular
    size, where legacy spectroscopic surveys and NED cross-identifications
    are most complete.}
  \label{fig:redshift-completeness}
\end{figure}

\subsection{Distances}
\label{sec:catalog:distances}

We assign the adopted distance (\texttt{DIST}) from four sources, in the
same decreasing-priority spirit as the redshift: high-precision
``gold-tier'' measurements compiled by NED from a single literature source,
LVD, NED's own blended mean distance, and, for every remaining galaxy with
a measured redshift, the Hubble-flow luminosity distance implied by
\texttt{Z\_COSMO} (Section~\ref{sec:catalog:redshifts}) under our fiducial
cosmology (Section~\ref{sec:intro}).  We adopt a gold-tier distance when
the highest-priority technique available for a galaxy---Cepheids, the
tip of the red-giant branch, the planetary-nebula luminosity function,
surface-brightness fluctuations, Type Ia or II supernovae, the
globular-cluster luminosity function, or SDSS-II Type Ia supernovae, in
that order, with typical distance-modulus uncertainties ranging from
0.08~mag for Cepheids to 0.26~mag for SDSS-II SNe~Ia---is confirmed by at
least three independent literature references; when it is not, but a
lower-priority technique is both well replicated (three or more references
agreeing to 15\%) and substantially discrepant from the top technique (by
more than 15\%), we adopt the better-replicated measurement instead,
preferring reproducibility over technique fidelity.  Failing a gold-tier
measurement we adopt the LVD distance, converted from kpc to Mpc; failing
that, NED's blended mean distance; and, for every remaining galaxy with a
measured redshift, the Hubble-flow distance implied by \texttt{Z\_COSMO},
so that essentially every galaxy with a redshift also carries a distance
estimate even absent a direct measurement.

We store the adopted distance, its inverse variance, its source
(\texttt{DIST}, \texttt{DIST\_IVAR}, \texttt{DIST\_REF} $\in
\{$\texttt{NED\_DIRECT}, \texttt{LVD}, \texttt{NED}, \texttt{ZCOSMO}$\}$),
and the specific technique behind it (\texttt{DIST\_METHOD}); as with
redshifts, \texttt{DIST\_IVAR} $= 0$ identifies a missing distance, and we
assign a 15\% fractional floor to any adopted distance with no reported
uncertainty rather than leave its inverse variance undefined.

\todo{Figure: compare the adopted distance against redshift-derived Hubble-flow
  distance (or against an independent distance compilation) as a function of
  \texttt{DIST\_REF}, to validate the priority scheme and quantify the
  Hubble-flow floor's accuracy for the bulk of the sample.}

%%----------------------------------------------------------------------
\section{Summary}
\label{sec:summary}
%%----------------------------------------------------------------------

We present the Siena Galaxy Atlas 2025 (SGA-2025), a multiwavelength
photometric catalog of $\sim\!486{,}000$ large, angularly resolved galaxies
selected from the final DESI Legacy Surveys Data Release 11 footprint.
The principal data products we deliver are:

\begin{enumerate}
  \item Postage-stamp image mosaics in up to ten bands ($grz/griz$,
    $W1$--$W4$, FUV, NUV) for each catalogued galaxy.
  \item Ellipse-fit geometry (position angle, axis ratio) and isophotal
    diameters at $\mu = 22$--26~mag~arcsec$^{-2}$.
  \item Curve-of-growth total magnitudes and half-light radii in all
    available bands.
  \item Spectroscopic redshifts, which we cross-match from LVD, DESI DR1,
    SDSS DR17, and NED.
  \item Adopted distances, drawn in priority order from NED gold-tier
    direct measurements, LVD, NED's blended mean distance, and, for the
    remainder, a Hubble-flow distance from the peculiar-velocity-corrected
    redshift.
\end{enumerate}

The SGA-2025 supersedes the SGA-2020 \citep{moustakas23a} and increases
the galaxy count by a factor of $\sim\!X$, extending coverage to bluer and
redder wavelengths.  Future work will incorporate improved photometric
redshifts, stellar mass estimates, and star-formation rate indicators.

%%----------------------------------------------------------------------
\begin{acknowledgments}

\textbf{Tom Jarrett} Dmitry, David Cooke and the NED
team. Rose. Trishotfury.

J.M.\ gratefully acknowledges support from XXX.

The DESI Legacy Imaging Surveys are supported by the U.S.\ Department of
Energy, Office of Science, Office of High Energy Physics, under Contract
No.\ DE-AC02-05CH11231, and by the National Energy Research Scientific
Computing Center, a DOE Office of Science User Facility, under the same
contract.

GALEX is a NASA Small Explorer mission.  The GALEX data used in this work
were obtained from the Mikulski Archive for Space Telescopes (MAST) at the
Space Telescope Science Institute.

This research has made use of the NASA/IPAC Extragalactic Database (NED),
which is funded by the National Aeronautics and Space Administration and
operated by the California Institute of Technology.

\end{acknowledgments}

\facilities{Blanco (DECam), Bok (BASS), Mayall (MzLS), WISE, GALEX}

\software{
  \texttt{astropy} \citep{astropy-collaboration22a},
  %\texttt{fitsio} \citep{sheldon23a},
  \texttt{healpy} \citep{zonca19a},
  \texttt{matplotlib} \citep{hunter07a},
  \texttt{numpy} \citep{harris20a},
  \texttt{scikit-learn} \citep{pedregosa11a},
  \texttt{scipy} \citep{virtanen20a},
  \texttt{seaborn} \citep{waskom21a},
  Tractor \citep{lang16a}
}

%%----------------------------------------------------------------------
\appendix

\section{Details of Parent-Sample and Group-Catalog Construction}
\label{sec:appendix:parent}

This appendix provides the full technical detail underlying the
parent-sample summary of Section~\ref{sec:parent:summary}: the six
external catalogs we merge into a single, uncut candidate list and the
cross-matching and geometry-standardization steps we apply to combine
them (Section~\ref{sec:input:parent:external}); the
visual-inspection-driven cleanup we apply to that no-cuts catalog
(Section~\ref{sec:parent:vicuts}); the self-supervised-learning
classification and photometric curation that produce the archive-stage
catalog (Section~\ref{sec:parent:archive}); the construction of the
Gaia+Tycho bright-star mask (Section~\ref{sec:parent:gaiamask}); and the
release-by-release, overlay-model assembly of the final parent sample
(Section~\ref{sec:parent:final}).

\subsection{External Catalogs}
\label{sec:input:parent:external}

Table~\ref{tab:parent-external} summarizes the six external catalogs we
merge to build the no-cuts parent sample, together with each catalog's
size prior to any cross-matching or cuts.  HyperLeda \citep{makarov14a}
is our primary input, since it is the most complete all-sky compilation
of galaxy positions and $D_{25}$ diameters available; we assembled it
from three successive database queries, executed on separate dates,
that together select confirmed galaxies, known multiple systems, and a
small number of remaining ambiguous-type entries -- the third query
was itself motivated by our discovery that nine SGA-2020 galaxies were
absent from the first two exports, all sharing an object-type code that
the initial query had inadvertently excluded.  We reconcile HyperLeda
against NED-LVS \citep[NED Local Volume Sample;][]{cook23a}, an
independently assembled NED cross-identification and redshift
compilation, to identify a high-confidence overlap sample and to recover
HyperLeda objects that NED resolves under a different name or position.
We supplement these two all-sky compilations with the Local Volume
Database \citep[LVD;][]{pace25a}, a hand-vetted census of confirmed
dwarf and Local Group galaxies within $\sim\!20$~Mpc that supplies the
most reliable positions, distances, and structural parameters for this
population; every SGA-2020 galaxy \citep{moustakas23a} within the DR11
footprint, carried forward directly; a hand-curated ``custom'' catalog
of one-off additions not resolved by any of the above, dominated by the
SMUDGes ultra-diffuse and low-surface-brightness galaxy sample
\citep{zaritsky23a}; and a supplemental list of large galaxies
identified during earlier (DR9/DR10) SGA processing but never folded
into a formal external catalog.

\begin{deluxetable*}{llrl}
\tablecaption{External Catalogs Merged into the SGA-2025 ``No-Cuts'' Parent Sample
  \label{tab:parent-external}}
\tablewidth{\textwidth}
\tablehead{\colhead{Catalog} & \colhead{Reference} & \colhead{$N$} & \colhead{Role}}
\startdata
HyperLeda            & \citet{makarov14a} & 4,855,956 & \parbox{3.0in}{Primary all-sky compilation of galaxy positions, diameters ($D_{25}$), axis ratios, position angles, and $B$-band magnitudes; combines separate ``galaxies,'' ``multiples,'' and ambiguous-type exports.} \\
NED-LVS               & \citet{cook23a}     & 1,872,475 & \parbox{3.0in}{NASA/IPAC Extragalactic Database Local Volume Sample; an independent cross-identification and redshift compilation against which we reconcile HyperLeda.} \\
LVD                    & \citet{leaman24a}   & 704       & \parbox{3.0in}{Local Volume Database; a hand-vetted compilation of confirmed dwarf and Local Group galaxies within $\sim\!20$~Mpc, with the most reliable positions, distances, and structural parameters for this population.} \\
SGA-2020               & \citet{moustakas23a}& 383,620   & \parbox{3.0in}{Prior SGA release; every SGA-2020 galaxy within the DR11 footprint is carried forward directly.} \\
Custom                 & \citet{zaritsky23a} & 2,927     & \parbox{3.0in}{Hand-curated one-off additions not resolved by any of the above catalogs, including the SMUDGes ultra-diffuse/low-surface-brightness galaxy sample.} \\
DR9/DR10 supplement    & \nodata             & 44,207    & \parbox{3.0in}{Large galaxies identified during earlier (DR9/DR10) SGA processing but not yet folded into a formal external catalog.} \\
\enddata
\tablecomments{$N$ is the number of rows in each cached external-catalog file prior to any cross-matching, deduplication, or angular-size cut; the same physical object is frequently present in more than one catalog (Section~\ref{sec:input:parent:merge}).}
\end{deluxetable*}


\subsubsection{Cross-Matching and Merging}
\label{sec:input:parent:merge}

We merge these six catalogs through a sequence of cross-matching stages,
each of which adds one catalog's remaining, not-yet-represented objects
to a running parent table.  We first deduplicate and clean the NED
cross-identification products for HyperLeda and NED-LVS individually --
dropping exact and near-position duplicates and known non-galaxy NED
object types (QSOs, stars, and \ion{H}{2} regions, among others) -- and
then join them against one another wherever NED independently resolves
both to the same object, which yields our highest-confidence
cross-identifications.  We next add whichever HyperLeda and NED-LVS
objects remain unmatched by this join, in turn, followed by any
remaining HyperLeda objects that NED does not resolve at all but that
have both a measured diameter and a plausible object classification.
Onto this NED/HyperLeda backbone we then merge LVD (first by NED name,
then by PGC number, then as new objects entirely), SGA-2020 (by PGC
number), the custom catalog (matching SMUDGes and other one-off entries
to existing rows within $5\arcsec$ before adding the remainder as new
objects; SMUDGes half-light radii are converted to Holmberg diameters
via a $1.2\times$ correction upstream of this step), and finally the
DR9/DR10 supplement.  Because no two literature catalogs agree
perfectly on object names, positions, or classifications, we hard-code
dozens of individual corrections throughout this process -- resolving
duplicate or conflicting NED cross-identifications, correcting
mis-assigned PGC numbers, and excluding known erroneous matches -- and
we guard every stage with internal consistency checks that raise an
error if an object's identifier is unexpectedly duplicated.

Once every catalog has been merged, we resolve each object's final
adopted position, name, and redshift by source priority: LVD takes
precedence for position and name, followed in order by NED, NED-LVS,
SGA-2020, HyperLeda, and the DR9/DR10 supplement, while for redshift we
instead prefer NED-LVS, then NED, then HyperLeda.  We apply an
analogous priority ordering -- LVD and the custom catalog ahead of
HyperLeda, NED-LVS, and the DR9/DR10 supplement -- to select each
object's adopted diameter, axis ratio, position angle, and magnitude,
while retaining separate, catalog-specific geometry columns (e.g.,
diameters from HyperLeda and from SGA-2020 individually) for later
comparison.

\subsubsection{Geometry Standardization}
\label{sec:input:parent:geometry}

Each external catalog reports size, shape, and magnitude information in
its own native system, so before the priority ordering described above
can be applied we first reduce every catalog to a common set of columns
-- diameter, axis ratio, position angle, and magnitude -- using
whichever conversion its native columns require.  For HyperLeda, we
adopt the $D_{25}$ diameter and axis ratio directly from its tabulated
$\log D_{25}$ and $\log R_{25}$ and pair them with its $B$-band total
magnitude.  For SGA-2020, we carry the diameter, axis ratio, position
angle, and magnitude forward unchanged.  For LVD, we convert the
tabulated half-light radius to a diameter via a $2.4\times$ correction
and take the axis ratio as one minus the reported ellipticity, pairing
both with the $V$-band apparent magnitude.  For the DR9/DR10
supplement, we derive the diameter, axis ratio, and position angle from
the Tractor S\'ersic shape parameters and adopt the $r$-band model flux
as the magnitude.  Where a catalog supplies no direct size or shape
measurement at all -- true of the raw NED cross-identification products
underlying the HyperLeda/NED-LVS join in
Section~\ref{sec:input:parent:merge} -- we instead merge together
whichever of the RC3, Two Micron All Sky Survey (2MASS), Sloan Digital
Sky Survey (SDSS), or ESO literature measurements NED provides, in that
decreasing order of priority.

\subsection{Visual-Inspection-Driven Cleanup}
\label{sec:parent:vicuts}

The no-cuts merge described above is deliberately inclusive: because
HyperLeda and NED resolve object names and positions independently and
imperfectly, the resulting candidate list also contains a substantial
population of non-galaxy contaminants -- literature galaxy-cluster and
-group catalogs, X-ray and radio source lists, deep-field photometric
catalogs, and other NED entries that are not resolved galaxies at all --
together with a smaller number of genuine cross-identification errors
and close, duplicate detections of the same object.  We remove this
contamination, and correct what visual inspection (VI) reveals to be
wrong, in a fixed sequence of cleanup steps applied to the no-cuts
catalog.  We first overwrite each affected object's coordinates and
geometry with hand-vetted values from a dedicated VI-properties file
(Table~\ref{tab:parent-vicuts}), which also corrects a further $\sim\!60$
SGA-2020 diameters we found, by inspection, to be grossly over- or
underestimated.  We then drop non-galaxy contaminants across the full
catalog using large, hand-curated lists of object-name prefixes and
embedded reference tags -- built up over many rounds of VI and
distinguishing contaminants that must always be dropped from those
protected whenever they carry a measured diameter -- and separately
correct NED cross-identification errors in which HyperLeda properties
(coordinates, diameter, PGC number) were erroneously associated with a
multi-galaxy \texttt{GPair}/\texttt{GTrpl} system entry rather than the
correct individual galaxy.  We next resolve close ($<1\arcsec$) pairs of
candidate entries that duplicate the same physical object, before
applying a final, narrower prefix-based cleanup that visually removes
remaining \texttt{GPair}/\texttt{GTrpl} merger systems: those lacking any
measured diameter are dropped outright, while those with a measured
diameter are dropped only if named explicitly in our VI-actions record.
Table~\ref{tab:parent-vicuts} summarizes the size of each hand-curated
resource this cleanup draws on; in total, the VI-actions record alone
encodes more than 32,000 individual object-level decisions.  Because
this cleanup is aggressive enough to risk discarding genuine
low-surface-brightness dwarfs, we verify after every step that every
LVD galaxy (Section~\ref{sec:input:parent:external}) remains present in
the catalog, and we abort the build if one has been inadvertently
dropped.  \todo{Report the net nocuts-to-vicuts object count for the
  final (v1.0) processing version once available; in the current
  development version this stage reduces the candidate list from
  $\sim\!5.10$ million to $\sim\!4.79$ million objects.}

\begin{deluxetable*}{llrl}
\tablecaption{Hand-Curated Resources Applied During Visual-Inspection-Driven Cleanup
  \label{tab:parent-vicuts}}
\tablewidth{\textwidth}
\tablehead{\colhead{Resource} & \colhead{Type} & \colhead{$N$} & \colhead{Role}}
\startdata
\texttt{SGA2025-vi-properties.csv}       & file        & 1,171  & \parbox{3.0in}{Hand-corrected coordinates and geometry (position, diameter, axis ratio, position angle) for individual objects identified during visual inspection.} \\
\texttt{SGA2025-crossid-errors.csv}      & file        & 1,777  & \parbox{3.0in}{NED cross-identification corrections: HyperLeda-derived properties mistakenly associated with a multi-galaxy \texttt{GPair}/\texttt{GTrpl} system entry instead of the correct individual galaxy.} \\
\texttt{SGA2025-vi-actions.csv}          & file        & 32,011 & \parbox{3.0in}{Object-level visual-inspection actions: 28,144 outright drops and 3,867 replacements of NED coordinates with HyperLeda coordinates.} \\
\texttt{LVD-geometry-updates-v1.0.5.csv} & file        & 165    & \parbox{3.0in}{Hand-corrected positions and structural parameters for individual LVD dwarf galaxies.} \\
Uncommon-prefix exclusion lists          & in-source   & 201    & \parbox{3.0in}{Hand-curated object-name prefixes (8 dropped only without a measured diameter, 119 dropped regardless) and embedded reference tags (74) identifying non-galaxy contaminants -- literature X-ray/radio source catalogs, deep-field photometric catalogs, and cluster/group member lists -- plus 14 individually vetoed object names always protected from removal.} \\
\enddata
\tablecomments{$N$ counts rows (for CSV files) or hand-curated list entries (for the in-source exclusion lists) as of the current processing version; see Section~\ref{sec:parent:vicuts}.}
\end{deluxetable*}


\subsection{SSL Classification and Photometric Curation}
\label{sec:parent:archive}

We apply a third round of automated and semi-automated curation to the
VI-cuts catalog, combining a machine-learning classification veto with
additional close-pair resolution, Local Group and Galactic-contaminant
flagging, and a photometry-based cut on the smallest candidates.  We
first remove every object flagged by a self-supervised-learning (SSL)
classifier trained to distinguish genuine galaxies from spurious or
non-galaxy sources, with three categories of exception: objects present
in a hand-curated veto list (Table~\ref{tab:parent-archive}), objects
already corrected by hand in the VI-properties file
(Section~\ref{sec:parent:vicuts}), and objects whose NED coordinates we
have already replaced with HyperLeda coordinates via the VI-actions
file; we treat any exception-list object absent from the classifier's
output as a consistency error, since these lists are meant to override
the classifier rather than introduce new objects.  \todo{The SSL
  classification step itself was run by a collaborator using an earlier,
  since-modified version of the ssl-legacysurvey pipeline (see
  \texttt{SGA/CLAUDE.md}); pin down the exact commit/methodology used
  and summarize it here, including how the training set was constructed
  and what performance we quote for it.}  We then resolve a second round
of close ($<5\arcsec$) pairs, breaking ties in favor of the lower-PGC
(and hence typically better-known) name; flag objects falling within the
Large or Small Magellanic Cloud's footprint or within a Galactic
globular cluster or planetary nebula's footprint, both of which we treat
separately from ordinary extragalactic sources; and remove small
candidates ($D_{\rm init} < 20\arcsec$) that fail an automated
photometry check, except those already protected by the custom,
properties, or VI-actions overrides, or those within a low-redshift
($z < 0.1$) projected radius of the Coma or Virgo cluster, whose
satellite populations we retain for dedicated visual inspection.
Finally, we add several boolean columns from additional hand-curated
action files (Table~\ref{tab:parent-archive}) -- flagging objects with
fixed ellipse-fitting geometry, relaxed or strengthened neighbor
masking, and Tractor sources forced to a point-source model -- and, for
every object, the projected distance to and masking radius of the
nearest bright Gaia star.  As in the previous stage, we verify after
every major removal step that no LVD dwarf has been inadvertently
dropped.  \todo{Report the net vicuts-to-archive object count for the
  final (v1.0) processing version once available; in the current
  development version this stage reduces the candidate list from
  $\sim\!4.79$ million to $\sim\!1.19$ million objects, the large drop
  being dominated by the $D_{\rm init} < 20\arcsec$ photometric cut.}

\begin{deluxetable*}{llrl}
\tablecaption{Hand-Curated Resources Applied During SSL Classification and Photometric Curation
  \label{tab:parent-archive}}
\tablewidth{\textwidth}
\tablehead{\colhead{Resource} & \colhead{Type} & \colhead{$N$} & \colhead{Role}}
\startdata
SSL veto lists                        & files (2) & 222 & \parbox{3.0in}{Objects manually exempted from the automated SSL-based classification veto (Section~\ref{sec:parent:archive}), on top of the properties- and VI-actions-file exemptions applied automatically.} \\
\texttt{SGA2025-resolved.csv}         & file      & 49  & \parbox{3.0in}{Objects flagged \texttt{RESOLVED}: too large or diffuse for automated ellipse-fitting, and therefore also forced \texttt{FIXGEO}.} \\
\texttt{SGA2025-fixgeo.csv}           & file      & 117 & \parbox{3.0in}{Objects for which we hold the elliptical geometry fixed during ellipse-fitting (Section~\ref{sec:pipeline:ellipse:mask}) rather than allowing moment-based refinement.} \\
\texttt{SGA2025-lessmasking.csv}      & file      & 23  & \parbox{3.0in}{Objects for which we relax the default masking of neighboring sources during ellipse-fitting.} \\
\texttt{SGA2025-moremasking.csv}      & file      & 15  & \parbox{3.0in}{Objects for which we apply more aggressive masking of neighboring sources during ellipse-fitting.} \\
\texttt{SGA2025-forcepsf.csv}         & file      & 9   & \parbox{3.0in}{Objects for which newly detected neighboring sources are forced to a point-source (PSF) Tractor model.} \\
\enddata
\tablecomments{$N$ counts rows (for CSV files) or object entries (for the SSL veto lists) as of the current processing version; see Section~\ref{sec:parent:archive}. Every listed action column (\texttt{IN\_RESOLVED}, \texttt{IN\_FIXGEO}, \texttt{IN\_LESSMASKING}, \texttt{IN\_MOREMASKING}, \texttt{IN\_FORCEPSF}) is carried forward into the final parent catalog and consulted directly by the ellipse-fitting pipeline (Section~\ref{sec:pipeline:ellipse}).}
\end{deluxetable*}


\subsection{Bright-Star Masking}
\label{sec:parent:gaiamask}

We build the Gaia+Tycho bright-star mask that underlies both the
\texttt{NEARSTAR}/\texttt{INSTAR} sample flags and the pixel-level
\texttt{brightstar}/\texttt{gaiastar} imaging-mask bits (Table~\ref{tab:maskbits-common})
in two stages: a one-time construction of an all-sky catalog of bright
and medium-bright stars and their adopted masking radii
(\texttt{SGA/archive/bin-SGA2025/build-gaia-mask}), and a per-catalog
association step that measures every parent-sample object's proximity
to its nearest such star (\texttt{SGA.parent.add\_gaia\_masking}).

To build the star-mask catalog, we loop over the full grid of DR9
imaging bricks in batches of 10,000, and for each brick query
\textsl{legacypipe}'s reference-source machinery
(\texttt{legacypipe.reference.get\_reference\_sources}) for every Gaia
DR3 and Tycho-2 star it flags as ``bright'' or ``medium,'' restricting
each per-brick query to a zero-margin footprint so that a star near a
brick boundary is picked up by exactly one neighboring brick rather
than double-counted.  We merge the resulting per-brick tables into a
single all-sky catalog and additionally write a Galactic-latitude-trimmed
($|b|>9\arcdeg$) ``extragalactic'' version, which is the one the
SGA-2025 pipeline actually consumes.

The adopted magnitude-radius relation --
\begin{equation}
  r_{\rm mask} = 1630\farcs0 \times 1.396^{-m_{\rm mask}},
  \label{eq:gaiamaskradius}
\end{equation}
where $m_{\rm mask}$ is the star's masking magnitude -- is the same
relation used by the DR11 \textsl{legacypipe} pipeline itself
(Section~\ref{sec:imaging:optical}) to mask reference sources during
Tractor photometry, so that the footprint we exclude from automated
galaxy detection is, by construction, consistent with the footprint
DR11 excludes from its own source modeling.  We define $m_{\rm mask}$
as the brighter of a star's optical magnitude (Gaia $G$, or, for
Tycho-2 stars, a Gaia-$G$-equivalent magnitude estimated from
$V_T$) and its $z$-band magnitude guess plus 1~mag, which expands the
mask appropriately around very red stars whose optical light
underrepresents their true angular extent.  We classify Gaia stars
with $m_{\rm mask}<13$ as ``bright'' and those with
$13\le m_{\rm mask}<16$ as ``medium,'' retaining both categories in the
mask catalog; Tycho-2 stars, which duplicate the brightest Gaia
sources, are instead restricted to $m_{\rm mask}<13$ only, dropping
roughly 13,000 of 2.5 million candidates (Table~\ref{tab:gaia-mask-params}).

Given this master catalog, \texttt{add\_gaia\_masking} measures, for
every object in a working parent-sample catalog, the angular
separation to its nearest bright or medium star, normalized by that
star's masking radius (\texttt{STARFDIST}); we process stars in
1-magnitude bins from faintest to brightest to keep each
nearest-neighbor search computationally tractable, and retain, for
each object, the smallest \texttt{STARFDIST} found across all bins
(capping consideration at twice a star's masking radius) so that the
final value reflects genuine proximity rather than an artifact of bin
ordering.  We call \texttt{add\_gaia\_masking} at two points in the
pipeline: once when the archive-stage catalog is finalized
(Section~\ref{sec:parent:archive}), and again for every subsequent
\texttt{build\_parent} release (Section~\ref{sec:parent:final}), since
each release's additions and geometry corrections can change which
star is nearest to a given object.  The resulting \texttt{STARFDIST}
values set the \texttt{NEARSTAR} and \texttt{INSTAR} \texttt{SAMPLE}
bits (Table~\ref{tab:sample-bitmask}) and, together with
\texttt{MCLOUDS} and \texttt{GCLPNE}, force the \texttt{NORADWEIGHT}
\texttt{ELLIPSEMODE} bit (Table~\ref{tab:ellipsemode-bitmask}) during
ellipse-fitting (Section~\ref{sec:pipeline:ellipse}), since
radially-weighted light centroids are unreliable for objects
contaminated by a nearby star's flux.

\begin{deluxetable*}{lcl}
\tablecaption{Gaia+Tycho Bright-Star Mask: Adopted Parameters
  \label{tab:gaia-mask-params}}
\tablewidth{\textwidth}
\tablehead{\colhead{Parameter} & \colhead{Value} & \colhead{Description}}
\startdata
$r_{\rm mask}$                    & $1630\farcs0 \times 1.396^{-m_{\rm mask}}$ & \parbox{3.6in}{Adopted magnitude-radius relation (Rongpu Zhou; identical to the one used by the DR11 \textsl{legacypipe} reference-source masking).} \\
$m_{\rm mask}$ (Gaia)              & $\min(G,\,z_{\rm guess}+1)$                & \parbox{3.6in}{Masking magnitude: the brighter of the Gaia $G$ magnitude and the $z$-band magnitude guess plus 1~mag, so masks expand appropriately around very red stars.} \\
Bright threshold (Gaia)            & $m_{\rm mask} < 13$                        & \parbox{3.6in}{Classified \texttt{isbright}.} \\
Medium threshold (Gaia)            & $13 \le m_{\rm mask} < 16$                 & \parbox{3.6in}{Classified \texttt{ismedium}; both bright and medium Gaia stars are retained in the mask catalog.} \\
Bright threshold (Tycho-2)         & $m_{\rm mask} < 13$                        & \parbox{3.6in}{Only bright Tycho-2 stars are retained (no medium category); drops $\sim\!13{,}000$ of 2.5 million Tycho-2 stars.} \\
Keep radius (Gaia)                 & $4\, r_{\rm mask}$                         & \parbox{3.6in}{Radius within which a source is retained as a Tractor reference source (e.g., for stellar-halo subtraction), not merely masked.} \\
Keep radius (Tycho-2)              & $2\, r_{\rm mask}$                         & \parbox{3.6in}{As above, for Tycho-2 stars.} \\
Galactic-latitude trim             & $|b| > 9\arcdeg$                           & \parbox{3.6in}{Footprint of the ``extragalactic'' mask catalog actually consumed by the SGA-2025 pipeline.} \\
\texttt{NEARSTAR} threshold        & $d/r_{\rm mask} < 1.2$                     & \parbox{3.6in}{\texttt{SAMPLE} bit (Table~\ref{tab:sample-bitmask}) set from \texttt{STARFDIST}.} \\
\texttt{INSTAR} threshold          & $d/r_{\rm mask} < 0.5$                     & \parbox{3.6in}{\texttt{SAMPLE} bit (Table~\ref{tab:sample-bitmask}) set from \texttt{STARFDIST}.} \\
\enddata
\tablecomments{$m_{\rm mask}$, $r_{\rm mask}$, and the bright/medium classification are computed once, for every star in the all-sky mask catalog, by \texttt{SGA/archive/bin-SGA2025/build-gaia-mask}; $d$ is the projected separation between a parent-catalog object and its nearest bright/medium star, computed by \texttt{SGA.parent.add\_gaia\_masking} (Section~\ref{sec:parent:gaiamask}).}
\end{deluxetable*}


\subsection{Assembling the Final Parent Sample}
\label{sec:parent:final}

We assemble the final parent sample not in a single pass, but across a
sequence of thirteen numbered releases (v0.30 through v1.6,
Table~\ref{tab:parent-overlays}), each of which refines the previous
release's \emph{fitted ellipse geometry} rather than rebuilding from the
archive-stage catalog described above.  Earlier releases, through
v0.22, were instead built directly from the archive catalog in a single
pass -- merging the per-region archive tables, appending hand-added
sources, applying visual-inspection corrections in place, and running
group-finding in one shot.  Beginning with v0.30 we adopted an
``overlay'' model in which every subsequent release is defined entirely
by a small, versioned set of hand-curated CSV files -- \texttt{adds.csv},
\texttt{drops.csv}, \texttt{updates.csv}, and \texttt{flags.csv} -- so
that the precise set of objects added, removed, corrected, or re-flagged
at each stage of visual inspection is preserved as a permanent,
inspectable record rather than folded irreversibly into the catalog.

For each release we first read and, where necessary, correct that
release's base ellipse catalogs across both imaging regions,
consolidating any object measured in both into a single row.  Several
early releases required release-specific restorations of this base
geometry -- for example, in the release immediately following the
transition to the overlay model, we restore the original (pre-fit)
position and geometry for every object that is an LVD dwarf, a member
of a multi-galaxy group, already flagged for a large Tractor-driven
centroid shift, or whose newly fitted diameter, axis ratio, or position
angle differs from its input value by more than 20\% -- a conservative
guard against a small fraction of erroneous fits before the pipeline
had matured.  We then apply that release's overlay files in order:
dropping listed objects (globally, or from a single imaging region);
appending hand-added objects, restoring their original identity
(\texttt{PGC}, \texttt{SGAID}) from the no-cuts catalog where a match
exists; and overwriting specific field values for individually
corrected objects.  We check the result for coincident (candidate
duplicate) detections and for unique object identifiers, repair each
object's region bits against the archive-stage catalogs, and
re-populate \texttt{SAMPLE} bits (\texttt{LVD}, \texttt{NEARSTAR}/
\texttt{INSTAR} from Gaia proximity, \texttt{MCLOUDS}, \texttt{GCLPNE})
and \texttt{ELLIPSEMODE} bits (from that release's consolidated
\texttt{flags.csv}) before deriving each object's ellipse-fitting
\texttt{FITMODE}.  We then re-run group-finding from scratch on the
full, corrected catalog (Section~\ref{sec:parent:groups}), excluding
\texttt{RESOLVED} and \texttt{FORCEPSF} objects (which we instead
assign to singleton groups), harmonize region bits within each
resulting group, and set the \texttt{OVERLAP} sample bit for every
group member whose (unscaled) ellipse overlaps that of another member
of the same group.  For the final several releases, we additionally
exempt any
galaxy group whose diameter did not change appreciably from
re-processing (via a per-release, opt-in list of previously converged
groups), so that the substantial computational cost of re-fitting a
group's mosaic (Section~\ref{sec:pipeline:ellipse}) is paid only for
groups that plausibly need it.

Because the overlay model records every visual-inspection decision as
an explicit, auditable file, we can quantify directly how much of this
work was done by hand.  Table~\ref{tab:parent-overlays} lists the
number of objects added, dropped, and individually corrected at each
release: in total, we made more than 45,000 individual VI-driven
additions, removals, or corrections across the v0.30--v1.6 releases,
dominated by two releases -- v0.40 and v1.2 -- that each followed a
particularly intensive round of visual inspection.  \texttt{v1.6} is
the final parent-sample version adopted for the SGA-2025 release.

\begin{deluxetable}{lcrrrr}
\tablecaption{Visual-Inspection Activity Recorded in the Per-Release Overlay Files
  \label{tab:parent-overlays}}
\tablehead{\colhead{Version} & \colhead{Base Version} & \colhead{Adds} & \colhead{Drops} & \colhead{Updates} & \colhead{Flags\tablenotemark{a}}}
\startdata
v0.30 & v0.22 &    27 &  4,024 &    417 & 291 \\
v0.40 & v0.30 &     7 & 12,679 &  1,177 & 324 \\
v0.50 & v0.40 &     0 &  3,468 &    194 & 328 \\
v0.60 & v0.50 &     0 &  2,275 &     78 & 334 \\
v0.70 & v0.60 &     0 &    153 &     32 & 340 \\
v0.80 & v0.70 &     0 &    114 &     43 & 338 \\
v1.0  & v0.80 &     1 &  1,227 &    528 & 345 \\
v1.1  & v1.0  &   257 &  1,603 &  7,198 & 348 \\
v1.2  & v1.1  & 5,082 &  2,602 &    265 & 350 \\
v1.3  & v1.2  &     5 &  1,318 &    292 & 417 \\
v1.4  & v1.3  &     2 &    273 &    450 & 458 \\
v1.5  & v1.4  &     0 &     29 &     48 & 467 \\
v1.6  & v1.4  &     0 &     31 &     54 & 468 \\
\hline
Total & \nodata & 5,381 & 29,796 & 10,776 & \nodata \\
\enddata
\tablenotetext{a}{The number of \texttt{ELLIPSEMODE} object-level bit assignments (\texttt{FIXGEO}, \texttt{RESOLVED}, etc.) recorded in that release's consolidated \texttt{flags.csv}; unlike Adds/Drops/Updates, this count is cumulative across releases, not a per-release increment, so it is not included in the total.}
\tablecomments{Each row corresponds to one call to \texttt{SGA.parent.build\_parent}, refining the ``Base Version'' release's ellipse-fitting output using the hand-curated \texttt{adds.csv}/\texttt{drops.csv}/\texttt{updates.csv}/\texttt{flags.csv} overlay files for ``Version'' (Section~\ref{sec:parent:final}); v0.22 is the final release built by the earlier, pre-overlay driver.  Versions v1.5 and v1.6 both use v1.4 as their base, since v1.6 -- the final SGA-2025 parent catalog -- was designed to reproduce v1.5 nearly exactly.}
\end{deluxetable}


%%----------------------------------------------------------------------

\bibliographystyle{aasjournalv7}
\bibliography{refs}

\end{document}
